\documentclass{article}
\usepackage[T1]{fontenc}
\usepackage{lmodern}
\usepackage{graphicx}
\usepackage{amsmath,amssymb}
\usepackage[a4paper,margin=2cm]{geometry}
\usepackage[absolute,overlay]{textpos}
\setlength{\TPHorizModule}{1mm}
\setlength{\TPVertModule}{1mm}
\usepackage[indonesian]{babel}
\usepackage{hyperref}
\setlength{\emergencystretch}{2em}
% unit: Halaman 72

\begin{document}

% === Halaman 72 (llm) ===

\begin{itemize}
  \item Serangan yang dilakukan terhadap A5/1 adalah \textit{known-plaintext attack}.
  \item Serangan semacam ini pertama kali dilakukan oleh Ross Anderson pada tahun 1994.
  \item Ross Anderson mencoba menerka 42 bit isi register R1 dan R2, dan menurunkan 23 bit R3 dari 42 bit tersebut. Pekerjaan menemukan kunci tersebut dilakukan pada komputer PC dan membutuhkan waktu komputasi lebih dari sebulan.
  \item Pada tahun-tahun selanjutnya, para kriptanalis berhasil menemukan kunci hanya dalam waktu beberapa menit saja.
  \item Secara umum, A5/1 sudah berhasil dikriptanalisis.
\end{itemize}

\end{document}
